\documentclass[UTF8]{ctexart}

\usepackage{geometry}
\geometry{margin=2.5cm}

\title{日志}
\author{王乐知 \and 徐晨曦}

\begin{document}

\maketitle

\noindent\textbf{2026年6月28日}

\vspace{1em}

在这里我们讨论论书的写作计划，日程安排等等。每次补充开一篇新的日志。

我们的名字叫什么好呢？还有大纲让我们来确认一下。

请使用简明英语语法手册.tex，大纲我弄好了。

直接在框架上修改，要先拉取再上传。

=======
\section*{推荐的协作流程}

注：该方案由AI友情提供

以后每次开始写文档时，建议按照这个顺序操作：

\begin{enumerate}
    \item \texttt{git pull} —— 获取同学最新修改。
    \item 编辑 \texttt{.tex} 文件。
    \item 本地编译检查是否正常。
    \item \texttt{git add .}
    \item \texttt{git commit -m "描述本次修改"}
    \item \texttt{git push}
\end{enumerate}

这样可以大幅减少冲突和覆盖彼此修改的风险。
>>>>>>> 6352a95 (Operate ways)

更新筛选文件，现在没有.pdf 文件了。

\end{document}